%
% Documento: Resumo (Português)
%

\begin{RESUMO}
\thispagestyle{empty}

    \noindent Este relatório apresenta o desenvolvimento de um chat distribuído
    executado por processos independentes e comunicado exclusivamente por sockets
    TCP. A solução permite configurar até 15 nós a partir de um catálogo estático
    de endereços, mantém relógios vetoriais e utiliza um nó coordenador como
    sequenciador das mensagens de grupo. Cada mensagem de grupo recebe um número
    de sequência global e é difundida por unicast a todos os participantes; em
    cada receptor, um buffer somente libera sequências contíguas. Também foram
    implementados detecção de falhas por batimentos periódicos e reeleição do
    coordenador pelo algoritmo do Valentão. Ensaios automatizados com sockets reais
    em uma única máquina, envolvendo 3, 8 e 15 processos independentes, confirmaram que
    todos os nós ativos entregaram as mensagens de grupo na mesma ordem e sem itens
    remanescentes no buffer. A análise do código e um ensaio de falha também
    identificaram limitações do protótipo: o recebimento de mensagens privadas, a
    captura de estado global e a continuidade do contador global após a troca de
    líder ainda não estão concluídos. Essas lacunas são documentadas para separar
    os requisitos demonstrados daqueles que dependem de evolução da implementação.

    \vspace*{0.5cm}\noindent\textbf{Palavras-chave}: sistemas distribuídos;
    comunicação de grupo; ordem total; relógio vetorial; eleição de líder.

\end{RESUMO}
